% !TeX root = ../slides.tex
% 01-intro.tex — motivation and problem statement.

\section{Motivation}

\begin{frame}{Minesweeper as a Sequential Decision Problem}
  \begin{highlightbox}{Standard RL Project goal}
    Implement and rigorously compare at least two distinct DRL techniques
    on a custom Minesweeper MDP.
  \end{highlightbox}
  \begin{itemize}
    \item Grid of $H \times W$ cells, $N$ hidden mines; revealed safe cells
          show the count of adjacent mines ($0$–$8$).
    \item Goal: uncover every safe cell without ever selecting a mine.
    \item Fully observable state (the player board fully summarizes what is
          known), but \alert{partial information} about mine locations drives
          the decision-making difficulty.
  \end{itemize}
\end{frame}

\begin{frame}{Anatomy of a Board}
  \centering
  \begin{tikzpicture}
    % Row 0 — all hidden
    \msc{0}{0}{mshidden}{}
    \msc{1}{0}{mshidden}{}
    \msc{2}{0}{mshidden}{}
    \msc{3}{0}{mshidden}{}
    % Row 1 — the clicked mine (red); one more mine hides at the right,
    % indistinguishable from an ordinary hidden cell
    \msc{0}{-1}{mshidden}{}
    \msc{1}{-1}{mshidden}{}
    \msc{2}{-1}{msmine}{\textcolor{white}{\faBomb}}
    \msc{3}{-1}{mshidden}{}
    % Row 2 — numbered border; the left cell hides a third mine
    \msc{0}{-2}{mshidden}{}
    \msc{1}{-2}{msrevealed}{\msnum{2}}
    \msc{2}{-2}{msrevealed}{\msnum{2}}
    \msc{3}{-2}{msrevealed}{\msnum{2}}
    % Row 3 — safely revealed
    \msc{0}{-3}{msrevealed}{\msnum{1}}
    \msc{1}{-3}{msrevealed}{\msnum{1}}
    \msc{2}{-3}{msrevealed}{}
    \msc{3}{-3}{msrevealed}{}
  \end{tikzpicture}

  \vspace{0.35cm}
  \mslegend{mshidden}{Hidden} \qquad
  \mslegend{msrevealed}{Revealed / number} \qquad
  \mslegend{msmine}{Mine that ended the episode}

  \vspace{0.25cm}
  \small A $4\times4$ example after a loss. Three mines sit on this board,
  but only the one just clicked is ever revealed — the numbers still
  correctly count \emph{all} neighboring mines, hidden or not.
\end{frame}

\begin{frame}{Why Minesweeper for RL}
  \twocol{
    \begin{itemize}
      \item Discrete, tunable state/action space ($H\!\times\!W$ scales
            difficulty).
      \item Sparse terminal reward (win/loss) plus dense shaping signals.
      \item Deterministic transition rule given the board $\Rightarrow$
            stochasticity comes only from mine placement.
      \item Existing DRL literature to compare against \parencite{wang2025minesweepercnn}.
    \end{itemize}
  }{
    \centering
    \begin{tikzpicture}[node distance=1.6cm]
      \node[state] (s) {$s_t$};
      \node[action, right=of s] (a) {$a_t$};
      \node[state, right=of a] (s2) {$s_{t+1}$};
      \draw[flowarrow] (s) -- (a);
      \draw[flowarrow] (a) -- (s2);
    \end{tikzpicture}
    \vspace{0.4cm}
    \small Uncover cell $\Rightarrow$ new board state, reward $r_t$.
  }
\end{frame}
